%  ----------
% | PACKAGES |
%  ----------

% INPUT, ENCODING & LANGUAGE
% ---
%\usepackage[utf8]{inputenc}          % XeLaTeX: non serve
\usepackage[english]{babel}
%\usepackage[T1]{fontenc}             % XeLaTeX: non serve

% -- Fonts per XeLaTeX (testo) --
% \usepackage[no-math]{fontspec}        % evita che fontspec tocchi i font matematici classici
% \defaultfontfeatures{Ligatures=TeX}
%\setmainfont{Latin Modern Roman}[
%  SmallCapsFont = Latin Modern Roman Caps,
%  SmallCapsFeatures = {Letters=SmallCaps}
% ]
% \setsansfont{Latin Modern Sans}
% \setmonofont{Latin Modern Mono}

% MATH & ALGORITHMS
% ---
\usepackage{amsmath}
\usepackage{amsfonts}
\usepackage{amssymb}
\usepackage{amsthm}
% \usepackage{algorithm2e}               % scegliamo questo come unico pacchetto per gli algoritmi
% \usepackage{algorithmic}             % <-- RIMOSSO per evitare doppioni/incompatibilità

% (Se preferisci lo stile algorithm + pseudocode, usa invece:)
\usepackage{algorithm}
\usepackage{algpseudocode} % da 'algorithmicx'
% % e in tal caso commenta \usepackage{algorithm2e} qui sopra.

% DOCUMENT STYLE & FORMATTING
% ---
\usepackage{style/isi_style_lt}
\usepackage{caption}
\usepackage{fancyhdr}
\usepackage{sectsty}
\usepackage{tocloft}
\usepackage{indentfirst}
\usepackage{microtype}

% COLORS, BOXES & HIGHLIGHTING
% ---
% \usepackage[usenames]{color}         % xcolor lo sostituisce
\usepackage[table]{xcolor}
\usepackage{tcolorbox}
\usepackage{framed}

% FIGURES, TABLES & LAYOUT
% ---
\usepackage{graphicx}
\usepackage{subcaption}                % sottofigure moderne
\usepackage{booktabs}
\usepackage{tabularx}
\usepackage{multirow}
\usepackage{arydshln}
\usepackage{adjustbox}
\usepackage{float}
\usepackage{array}

% CODE & LISTINGS
% ---
\usepackage{listings}
\usepackage{fancyvrb}

% SYMBOLS & ICONS
% ---
\usepackage{marvosym}
\usepackage{pifont}
\usepackage{arydshln}

% SPECIAL STRUCTURES
% ---
\usepackage{multicol}
\usepackage{dirtree}
\usepackage{tikz}
\usepackage{pgfplots}
\usepackage{svg}

%  -----------------
% |  CUSTOM COLORS  |
%  -----------------

% Rounded tag
\definecolor{tagbackground}{HTML}{EBEBEB}
\definecolor{tagborder}{HTML}{888888}
\definecolor{tagforeground}{HTML}{000000}

% Tables
\definecolor{tablehighlight}{HTML}{F3F3F3}
\definecolor{tablebordercolor}{HTML}{363636}
\definecolor{cellhighlight}{HTML}{E9E9E9}

% Figures
\definecolor{plotbackground}{HTML}{F3F3F3}
\definecolor{gray}{HTML}{394144}
\definecolor{blue-1}{HTML}{1876D2} % dark
\definecolor{blue-2}{HTML}{60D0E2}
\definecolor{orange-1}{HTML}{FBA418} % dark
\definecolor{orange-2}{HTML}{FDCC7B}
\definecolor{purple-1}{HTML}{4D4DFF} % dark
\definecolor{purple-2}{HTML}{8F8FFF}
\definecolor{red-1}{HTML}{F00000} % dark
\definecolor{red-2}{HTML}{FF7575}

% Output
\definecolor{baseline}{HTML}{F2F2F2}
\definecolor{gom}{HTML}{F0F4FF}
\definecolor{error}{HTML}{D10000}
\definecolor{correct}{HTML}{0A9E4D}

\pgfdeclareverticalshading{bluegradient}{100bp}{
  color(0bp)=(blue-1);
  color(30bp)=(blue-1);
  color(60bp)=(blue-2);
  color(100bp)=(blue-2)
}
\pgfdeclareverticalshading{orangegradient}{100bp}{
  color(0bp)=(orange-1);
  color(30bp)=(orange-1);
  color(60bp)=(orange-2);
  color(100bp)=(orange-2)
}
\pgfdeclareverticalshading{purplegradient}{100bp}{
  color(0bp)=(purple-1);
  color(30bp)=(purple-1);
  color(60bp)=(purple-2);
  color(100bp)=(purple-2)
}
\pgfdeclareverticalshading{redgradient}{100bp}{
  color(0bp)=(red-1);
  color(30bp)=(red-1);
  color(60bp)=(red-2);
  color(100bp)=(red-2)
}

% REFERENCES & HYPERLINKS
% ---
\usepackage{csquotes}                  % Raccomandato da biblatex per le citazioni
\usepackage{url}
\usepackage{hyperref}
% (Opzionale ma consigliato per gli indici PDF più stabili)
\usepackage{bookmark}

% BIBLIOGRAPHY
% ---
\usepackage[backend=biber,style=numeric,sorting=none,maxbibnames=99]{biblatex}
\addbibresource{biblio.bib}      % file .bib principale